\section{The remaining three-primary growth gate has rank one}
\label{sec:cubic-rank-gate}

The finite logarithmic certificate has a strong consequence when
combined with the rational complex-conjugation symmetry of this
particular tower.  We use $X'$, $T$, $r$, $G$ and $\gamma$ from
AC-K16, with $\varepsilon_3(\gamma)=4$.

\begin{theorem}[AC-K17: rank alternatives and a finite arithmetic test]
\label{AC-K17}
Let $r=\operatorname{rank}_{\mathcal O}X'$, so $0\leq r\leq7$.
If $r=0$, the sequence $(m_e)_{e\geq1}$ is uniformly bounded.
If $r\geq2$, then
\[
                         m_e=4\quad(e\geq1).           \tag{KR.1}
\]
If $r=1$, there are $g\in1+3\mathbb Z_3$ with
$v_3(g-1)=1$ and $k\in\{5,6\}$ such that
\[
 T\simeq\mathbb F_3[B]/(B^k),\quad B=\gamma-1,
 \qquad
 m_e=2+2v_3(4^eg-1)\quad(e\geq1).                     \tag{KR.2}
\]
In this last case the finite weight-zero invariant of the
first cyclotomic layer necessarily satisfies
\[
 \#\widetilde{\operatorname{Cl}}_3(E_1)=3^7,\qquad
 E_1=\mathbb Q(\mu_9,\sqrt[3]{269}).                    \tag{KR.3}
\]
Consequently an unconditional upper bound strictly less than
$3^7$ for that logarithmic class group would prove that
$(m_e)$ is uniformly bounded.  The tentative arithmetic output
for $E_1$ is not used in this theorem.
\end{theorem}
\begin{proof}
\proofstep{1. The rational characteristic polynomial.}
Complex conjugation $c$ acts on the whole field $E_\infty$.
It conjugates $\sigma$ to $\sigma^{-1}$ and commutes with
$\gamma$.  These identities follow from
\[
 \operatorname{Gal}(E_\infty/\mathbb Q)
       \simeq S_3\times\mathbb Z_3:
\]
the fiber product of the $S_3$-extension $E/\mathbb Q$
and the abelian cyclotomic extension over their common quadratic
field identifies the cyclotomic pro-three factor as central.
The induced action of $c$ on $X'$ is therefore semilinear for
$\omega\mapsto\omega^2$, and commutes with $\gamma$.
It preserves the torsion module $T$ and hence acts on the
free quotient $\mathcal O^r$.
If this semilinear map has matrix $H$, then commutation gives
$G=H\overline G H^{-1}$.  Thus
\[
 P(z)=\det(zI_r-G)\in\mathbb Z_3[z].                   \tag{KR.4}
\]
The finite image of the pro-three group generated by $\gamma$
on the free quotient modulo $\lambda$ is a three-group.
Every eigenvalue of this image is one in characteristic three:
if $M^{3^s}=I$, then $(M-I)^{3^s}=0$.
Consequently
\[
                    P(z)\equiv(z-1)^r\pmod3,          \tag{KR.5}
\]
since reduction of coefficients in $\mathbb Z_3$ modulo
$\lambda$ is the same as reduction modulo three.

The certified base logarithmic group in AC-K13 has order
$27$.  Applying the snake-lemma proof of (KI.7) also at
$e=0$ is legitimate: $X'/(\gamma-1)X'$ is this finite
logarithmic group, so the induced endomorphism $G-I_r$
on the free quotient is injective with finite cokernel.
It gives
\[
 3=t_0+2v_3(P(1)),\qquad
 t_0=\operatorname{length}_{\mathcal O}T/(\gamma-1)T.
                                                               \tag{KR.6}
\]
For $r>0$, (KR.5) implies $v_3(P(1))\geq1$.
The integer $t_0$ is nonnegative, so (KR.6) forces
\[
                     v_3(P(1))=1,\qquad t_0=1.        \tag{KR.7}
\]
In particular $T$ is then nonzero.

\proofstep{2. The finite torsion is killed by $\lambda$.}
Assume $r>0$ and write $B=\gamma-1$ on $T$.
Since $T$ is finite and the cyclotomic group is pro-three,
$B$ is nilpotent modulo three; because a power of three
annihilates $T$, some power of $B$ annihilates $T$ as well.
The finite module $T$ is thus a module over the local ring
$\mathbb Z_3[[B]]$ with maximal ideal $(3,B)$.
By (KR.7), $T/BT$ has order three, and
$T/(3,B)T=\mathbb F_3$.  Nakayama's lemma shows that $T$
is cyclic over this ring.

Conjugation $c$ commutes with $B$ and with $\mathbb Z_3$.
It acts by one sign $\epsilon\in\{1,-1\}$ on $T/(3,B)T$.
If $v$ is a generator, then
$v_\epsilon=(v+\epsilon c(v))/2$ has the same nonzero image
in this quotient.  It is again a generator and satisfies
$c(v_\epsilon)=\epsilon v_\epsilon$.
Since $c$ commutes with every coefficient and with $B$,
it acts by that scalar on all of $T$.
On the other hand $c\sigma c^{-1}=\sigma^{-1}$.
A scalar $c$ commutes with $\sigma$, so
$\sigma=\sigma^{-1}$ on $T$; together with $\sigma^3=1$
this gives $\sigma=1$.  Hence $\lambda T=0$.
Cyclicity and nilpotence now give
\[
                 T\simeq\mathbb F_3[B]/(B^k)
                         \quad\text{for some }k\geq1. \tag{KR.8}
\]
Since $4^e=1$ on this characteristic-three module, its
coinvariant length is exactly one for every twist $e$.
Moreover the $\mathcal O$-module sequence
$0\to T\to X'\to\mathcal O^r\to0$ splits as modules,
so the minimal number of generators of $X'$ is $r+k$.
AC-K16 therefore gives
\[
                              r+k\in\{6,7\}.           \tag{KR.9}
\]

\proofstep{3. All free ranks at least two.}
Suppose $r\geq2$.  Equation (KR.5) gives $P'(1)\equiv0\pmod3$.
For every $z\in1+3\mathbb Z_3$, polynomial Taylor expansion
with integral coefficients gives
\[
                         P(z)\equiv P(1)\pmod9.
\]
Indeed its linear term is a product of two multiples of three,
and its terms of degree at least two contain $(z-1)^2$.
By (KR.7), $v_3(P(z))=1$ throughout that set.
Taking $z=4^{-e}$, the determinant in (KI.6) differs from
$P(4^{-e})$ only by a unit.  Its $\lambda$-valuation is two;
the torsion contribution is one by (KR.8).  Thus $a_e=3$
and $m_e=4$, proving (KR.1).

\proofstep{4. The single remaining free rank.}
If $r=1$, equation (KR.4) says that $G$ is the scalar
$g\in\mathbb Z_3^\times$.  Equation (KR.5) gives
$g\in1+3\mathbb Z_3$, and (KR.7) gives $v_3(g-1)=1$.
Equations (KI.6) and (KR.8) then yield (KR.2).
The determinant cannot vanish at any positive integer $e$,
by the already proved finiteness of $W_e$.
Equation (KR.9) says that $k=5$ or $6$.

The cyclotomic extension above $E_1$ has generator $\gamma^3$
and the same split unramified module $X'$.
Thus its logarithmic class group is
$X'/(\gamma^3-1)X'$, by (KL.9) applied over $E_1$.
On the finite module (KR.8),
$(1+B)^3-1=B^3$, so the cokernel has length
$\min(k,3)=3$.  On the free rank-one quotient, lifting the
exponent gives $v_3(g^3-1)=2$, hence cokernel length four
over $\mathcal O$.  The same exact snake-lemma sequence
as (KI.7) therefore proves total length seven.
This proves (KR.3), including finiteness in this case.

Finally, if $r=0$, $X'$ is itself finite.  Every $W_e$ is
a quotient of that fixed module, so
$m_e\leq1+\operatorname{length}_{\mathcal O}X'$ uniformly.
The final implication follows by eliminating $r=1$ with
(KR.3), then applying either this bound or (KR.1).
\end{proof}

\paragraph{What this does and does not settle.}
The only free-rank case that can still carry unbounded
three-primary lengths is the single scalar specialization
in (KR.2).  Writing $g=4^{-\beta}$ with $\beta\in\mathbb Z_3$
would rewrite its valuation in terms of $e-\beta$; finiteness
at every positive integer does not bound that approximation.
The finite logarithmic test (KR.3) is therefore a strictly
smaller arithmetic gate than an unspecified infinite list
of exponents.  No uncertified class-group output is a proof
of its failure.  Even a uniform fixed-field three-primary
bound would still leave the other coefficient primes and
the source's growing CM hosts in H04.4.

In geometric language, conjugation acts on the compact
abelianized covering module from SGA~1's Galois category;
it is semilinear over the integral cubic character order.
The determinant (KR.4) is the characteristic polynomial
of that arithmetic module, not a motive's archimedean
zero operator.  The argument uses the exact K-theoretic
and logarithmic-class comparisons already cited in
AC-K10 and AC-K13--K14, and asserts no new scan-level
attribution in Grothendieck's archive.
